\section*{Hoofdstuk \ref{ch:g10:earth-energy-budget} --- De energiebalans van de Aarde}

\begin{solution}{exo:g10:earth-energy-budget:1}
De \omterm{def:g10:earth-energy-budget:insolation}{instraling} van een plaats is het vermogen van het zonlicht dat een
horizontale vierkante meter aan de top van de atmosfeer erboven ontvangt
(gemiddeld over een dag of een jaar). Het \omterm{def:g10:earth-energy-budget:albedo}{albedo} van een oppervlak is het deel
van het zonlicht dat erop valt en dat het weerkaatst. Verse sneeuw (0.8 tot 0.9)
weerkaatst veel meer dan de open oceaan (0.06).
\end{solution}

\begin{solution}{exo:g10:earth-energy-budget:2}
Weerkaatst: $100/340 \approx \qty{29}{\percent}$; geabsorbeerd door de lucht:
$80/340 \approx \qty{24}{\percent}$; geabsorbeerd door het oppervlak: $160/340
\approx \qty{47}{\percent}$.
\end{solution}

\begin{solution}{exo:g10:earth-energy-budget:3}
Asfalt: $800 \times (1 - 0.04) = \qty{768}{W}$. Sneeuw: $800 \times (1 -
0.85) = \qty{120}{W}$, meer dan zes keer minder.
\end{solution}

\begin{solution}{exo:g10:earth-energy-budget:4}
Verdamping: \qty{82}{W/m^2}, tegen \qty{21}{W/m^2} voor het opwarmen van de
lucht. Het water dat verdampt, onttrekt zijn \omterm{def:g10:earth-energy-budget:latent-heat}{latente warmte} aan de zee; de damp
stijgt met de lucht op, en waar hij condenseert tot wolkendruppeltjes, vaak
kilometers hoger, geeft hij die energie terug aan de lucht eromheen.
\end{solution}

\begin{solution}{exo:g10:earth-energy-budget:5}
In december helt het noordelijk halfrond van de Zon af: op de pool blijft de
Zon de hele dag onder de horizon, dus er komt geen zonlicht aan. Bij de
decemberzonnewende geldt dat overal voorbij $90^\circ - 23.44^\circ =
\qty{66.56}{\degree}$ NB.
\end{solution}

\begin{solution}{exo:g10:earth-energy-budget:6}
$1361 \sin 90^\circ = \qty{1361}{W}$; $1361 \sin 30^\circ =
\qty{680.5}{W}$; $1361 \sin 10^\circ \approx \qty{236}{W}$. De helft van $S_0$
vraagt $\sin h = 0.5$, dus $h = \qty{30}{\degree}$.
\end{solution}

\begin{solution}{exo:g10:earth-energy-budget:7}
Weerkaatst 100, geabsorbeerd door de lucht 80, geabsorbeerd door het oppervlak
160: $100 + 80 + 160 = 340$, al het zonlicht is verantwoord. Aan het oppervlak
geldt $185 = 25 + 160$: wat niet wordt weerkaatst, wordt geabsorbeerd.
\end{solution}

\begin{solution}{exo:g10:earth-energy-budget:8}
Schijf: $\pi \times (6.371\times10^6)^2 \approx \qty{1.275e14}{m^2}$; vermogen:
$1361 \times 1.275\times10^{14} \approx \qty{1.74e17}{W}$. Bol: $4\pi
R^2 \approx \qty{5.10e14}{m^2}$. Gemiddeld: $1.74\times10^{17}/5.10\times
10^{14} \approx \qty{340}{W/m^2}$, een kwart van $S_0$.
\end{solution}

\begin{solution}{exo:g10:earth-energy-budget:9}
$0.7 \times 5.10\times10^{14} \times 3.16\times10^7 \approx
\qty{1.1e22}{J}$ per jaar.
\end{solution}

\begin{solution}{exo:g10:earth-energy-budget:10}
De wolken weerkaatsen $100 - 53 = \qty{47}{W/m^2}$ meer zonlicht, dat dus niet
wordt geabsorbeerd. Ze houden $267 - 239 = \qty{28}{W/m^2}$ infrarood tegen.
Per saldo verliest de Aarde $47 - 28 = \qty{19}{W/m^2}$: de wolken koelen haar
als geheel af.
\end{solution}

\begin{solution}{exo:g10:earth-energy-budget:11}
Van ongeveer 10 mei tot eind juli, dus in mei, juni en juli. Op de evenaar
staat de Zon bij de twee equinoxen, in maart en september, om 12 uur 's middags
in het zenit; bij de zonnewenden passeert ze \qty{23.44}{\degree} van het
zenit, zodat de evenaar in juni en in december iets minder ontvangt.
\end{solution}

\begin{solution}{exo:g10:earth-energy-budget:12}
Per seconde: $82/(2.47\times10^6) \approx \qty{3.32e-5}{kg}$. Per dag:
$3.32\times10^{-5} \times 86\,400 \approx \qty{2.9}{kg}$, een laag van ongeveer
\qty{2.9}{mm}. Per jaar: $3.32\times10^{-5} \times 3.16\times10^7 \approx
\qty{1050}{kg}$, ongeveer \qty{1.05}{m} water.
\end{solution}

\begin{solution}{exo:g10:earth-energy-budget:13}
$S_0 \sin\delta = S_0/\pi$ geeft $\sin\delta = 1/\pi \approx 0.318$, dus
$\delta \approx \qty{18.6}{\degree}$. Dan is $\sin\lambda = 0.318/\sin 23.44^\circ
\approx 0.318/0.398 \approx 0.800$: $\lambda \approx \qty{53.15}{\degree}$ of
$180^\circ - 53.15^\circ = \qty{126.85}{\degree}$. Daartussen neemt $\lambda$
met \qty{73.7}{\degree} toe, dat is $73.7/360$ van een jaar, ongeveer
\num{75} dagen, van half mei tot eind juli.
\end{solution}

\begin{solution}{exo:g10:earth-energy-budget:14}
Het ijs absorbeert $1 - 0.6 = 0.40$ van het zonlicht, de oceaan $1 - 0.06 =
0.94$. Extra vermogen: $(0.94 - 0.40) \times 300 \times 10^{12} \approx
\qty{1.6e14}{W}$, meer dan alle fotosynthese van de planeet. De extra energie
warmt het water op, dat aan de randen van de ijsschotsen meer ijs doet smelten,
waardoor meer donker water bloot komt te liggen: de verandering voedt zichzelf.
\end{solution}

\begin{solution}{exo:g10:earth-energy-budget:15}
Dichtstbij: $1361 \times (149.6/147.1)^2 \approx \qty{1408}{W/m^2}$; verst:
$1361 \times (149.6/152.1)^2 \approx \qty{1317}{W/m^2}$; ongeveer
\qty{7}{\percent} verschil. Op \qty{60}{\degree} NB is de \omterm{def:g10:earth-energy-budget:insolation}{instraling} in juni
twintig keer groter dan in december: een verandering van een paar procent kan
dat niet veroorzaken. Bovendien staat de Aarde begin januari het dichtst bij
de Zon, tijdens de noordelijke winter, en hebben de twee halfronden, op dezelfde
afstand van de Zon, tegengestelde seizoenen: de helling van de as, niet de
afstand, maakt de seizoenen.
\end{solution}

\begin{solution}{pb:g10:earth-energy-budget:1}
\textbf{1.} $\pi \times (6.371\times10^6)^2 \approx \qty{1.275e14}{m^2}$.

\textbf{2.} $1361 \times 1.275\times10^{14} \approx \qty{1.735e17}{W}$.

\textbf{3.} $4\pi R^2 \approx \qty{5.101e14}{m^2}$.

\textbf{4.} $1.735\times10^{17}/5.101\times10^{14} \approx \qty{340}{W/m^2}$,
dat is $S_0/4$.

\textbf{5.} Evenaar: \qty{1361}{W}; \qty{60}{\degree} NB: $1361 \sin
30^\circ = \qty{680.5}{W}$.

\textbf{6.} $100/340 \approx 0.294$.

\textbf{7.} $340 - 100 = \qty{240}{W/m^2}$; over de hele Aarde $240 \times
5.101\times10^{14} \approx \qty{1.224e17}{W}$.

\textbf{8.} $53/340 \approx 0.156$. Wolken verhogen het \omterm{def:g10:earth-energy-budget:albedo}{albedo} met een factor
$100/53 \approx 1.9$: ze verdubbelen het bijna.

\textbf{9.} $(240/5.67\times10^{-8})^{1/4} \approx \qty{255}{K}$, ongeveer
\qty{33}{K} onder de \qty{288}{K} van het oppervlak.

\textbf{10.} $160 + 342 = \qty{502}{W/m^2}$.

\textbf{11.} $82 + 21 + 398 = \qty{501}{W/m^2}$. Het oppervlak wint ongeveer
\qty{1}{W/m^2}: de afgeronde onbalans van \qty{0.7}{W/m^2}.

\textbf{12.} $342/160 \approx 2.1$. De lucht stuurt het oppervlak als infrarood
meer dan twee keer zoveel energie als het zonlicht dat het absorbeert: het
broeikaseffect.

\textbf{13.} $82 \times 3.156\times10^7/(2.47\times10^6) \approx
\qty{1050}{kg}$, een laag van ongeveer \qty{1.05}{m} dik.

\textbf{14.} $0.7 \times 5.101\times10^{14} \approx \qty{3.6e14}{W}$.

\textbf{15.} $2805/72.07 \approx \qty{38.9}{kJ}$ per gram koolstof.

\textbf{16.} $104.9\times10^9$ ton is $1.049\times10^{17}$~g;
$1.049\times10^{17} \times 38.9\times10^3 \approx \qty{4.08e21}{J}$.

\textbf{17.} $4.08\times10^{21}/3.156\times10^7 \approx \qty{1.29e14}{W}$.

\textbf{18.} $1.29\times10^{14}/5.101\times10^{14} \approx
\qty{0.25}{W/m^2}$.

\textbf{19.} $0.7/0.25 \approx 2.8$: de onbalans van de Aarde is bijna drie
keer het vermogen dat fotosynthese opslaat.

\textbf{20.} $1.29\times10^{14}/1.224\times10^{17} \approx 0.0011$:
fotosynthese slaat \textbf{ongeveer \qty{0.1}{\percent}} op van het zonlicht
dat de Aarde absorbeert --- een duizendste.
\end{solution}
